\section*{LỜI NÓI ĐẦU}
\thispagestyle{empty}
Sự phát triển nhanh chóng của các thiết bị IoT ứng dụng cho các thiết bị đeo tay, xe tự hành hay một mạng lưới các thiết bị thông minh ngày càng xuất hiện nhiều hơn. Đi cùng với nó là nguy cơ bị đánh cắp dữ liệu, rò rỉ thông tin và mất quyền truy cập quản trị bởi những kẻ tấn công công nghệ cao. Việc xây dựng một công nghệ bảo mật thông tin, chẳng hạn như xác thực phần cứng hoặc chữ kí số là rất cần thiết. Nguyên lý Kerckhoffs nói rằng một hệ thống mật mã phải an toàn ngay cả khi kẻ tấn công biết được thuật toán mã hóa và giải mã, ngoại trừ khóa bí mật. Do đó, khóa ngẫu nhiên bí mật là một yếu tố quan trọng trong giao thức xác thực. Trong nỗ lực giải quyết vấn đề trên, Physical Unclonable Function (PUF) là công nghệ đầy hứa hẹn trong bảo mật phần cứng các thiết bị IoT. 

Phần này trình bày một cách rất khái quát (khoảng 1-2 trang) về bối cảnh hình thành và mục đích của đồ án. Lời cảm ơn với những tổ chức và cá nhân góp phần trong việc hoàn thiện đồ án nên đặt ở cuối mục này.
\cleardoublepage